\section{Method}
\label{sec:method}

Our semantic correspondence pipeline relies on dense features extracted from pretrained visual foundation models. In this section, we describe
the problem formulation, the training-free baseline based on cosine similarity, the light fine-tuning strategy adopted for different backbones,
and the prediction refinement using windowed soft-argmax.

%-------------------------------------------------------------------------

\subsection{Problem Setup}

Given a source image $I_s$ and a target image $I_t$, the goal of semantic correspondence is to predict, for each annotated keypoint $p_s$ in $I_s$,
the location $\hat{p}_t$ in $I_t$ corresponding to the same semantic object part. We evaluate our method on datasets providing sparse keypoint
annotations, namely SPair-71k~\cite{min2019spair} and PF-Willow, using the Percentage of Correct Keypoints (PCK) metric at multiple thresholds.

%-------------------------------------------------------------------------

\subsection{Training-Free Correspondence Baseline}

We first establish a training-free baseline following prior work on emergent correspondence from pretrained models~\cite{tang2023emergent, zhang2023tale}.
Given an image $I$, we extract a dense feature map $F(I) \in \mathbb{R}^{H' \times W' \times D}$ using a frozen backbone. We evaluate three pretrained
models: DINOv2~\cite{oquab2023dinov2}, DINOv3~\cite{simeoni2025dinov3}, and Segment Anything (SAM)~\cite{kirillov2023segment}. Features are extracted online without caching.

For each source keypoint $p_s$, we extract the corresponding feature vector $f_s = F(I_s)[p_s]$ and compute cosine similarity with all spatial locations
in the target feature map $F(I_t)$. The predicted correspondence is obtained by selecting the target location that maximizes the similarity score:

\[
\hat{p}_t = \arg\max_{r} \ \cos\big(f_s, F(I_t)[r]\big),
\]

where $r$ indexes spatial positions in the target feature map. This simple global argmax matching serves as a strong baseline for semantic correspondence.

%-------------------------------------------------------------------------

\subsection{Light Fine-Tuning of Pretrained Backbones}

While the training-free setting already provides strong correspondence performance, we further investigate light fine-tuning to adapt pretrained representations
to the semantic correspondence task. Fine-tuning is performed using keypoint supervision from SPair-71k by unfreezing only a limited subset of layers
for each backbone, while keeping the remaining parameters frozen.

We experiment with unfreezing the last few transformer blocks together with the final normalization layer. In particular, for DINO models we report results
obtained by unfreezing the last two transformer layer and the normalization layer. For SAM, we fine-tune the last two layers and the neck component
while keeping the remaining layers of the encoder frozen.
This design allows the model to adapt to correspondence-specific cues while preserving most of the general semantic structure learned during pretraining.

Training is supervised using a contrastive loss defined on feature representations at corresponding and non-corresponding locations. Given a source
keypoint $p_s$ with feature vector $f_s$ and its ground-truth correspondence $p_t$ in the target image with feature $f_t$, the objective encourages
high similarity between matching features while discouraging similarity with non-matching target locations. Specifically, for a batch of $N$ keypoints,
we optimize a contrastive objective of the form:

\[
\mathcal{L} = -\frac{1}{N} \sum_{i=1}^{N}
\log \frac{\exp(\cos(f_s^i, f_t^i) / \tau)}
{\sum_{j=1}^{N} \exp(\cos(f_s^i, f_t^j) / \tau)},
\]

where $\tau$ is a temperature hyperparameter. This formulation encourages discriminative feature alignment across image pairs while remaining robust
to intra-class variability and geometric differences.

%-------------------------------------------------------------------------

\subsection{Prediction Refinement with Windowed Soft-Argmax}

The global argmax operation produces discrete predictions and can be sensitive to local noise in the similarity map. To improve robustness and enable
sub-pixel localization, we adopt a windowed soft-argmax refinement strategy inspired by~\cite{zhang2024telling}.

For each source keypoint, we first obtain a coarse correspondence $\hat{p}_t$ using global argmax. We then define a local square window centered
at $\hat{p}_t$ with radius $r$. Within this window, we apply a softmax over similarity scores and compute a weighted average of spatial coordinates:

\[
\tilde{p}_t = \sum_{r \in \mathcal{W}(\hat{p}_t)} r \cdot 
\frac{\exp(\alpha \cdot S(r))}{\sum_{k \in \mathcal{W}(\hat{p}_t)} \exp(\alpha \cdot S(k))},
\]

where $S(r)$ denotes the cosine similarity at position $r$, $\mathcal{W}(\hat{p}_t)$ is the local window, and $\alpha$ controls the sharpness of
the distribution. We empirically observe that large window radii (\eg $r=10$) degrade performance, and therefore use smaller values
selected through validation.

%-------------------------------------------------------------------------

\subsection{Implementation Details}

Images are resized according to the backbone requirements: $518 \times 518$ for DINOv2, $512 \times 512$ for DINOv3 and for SAM.
For DINO-based models, we apply ImageNet normalization with mean $(0.485, 0.456, 0.406)$ and standard deviation $(0.229, 0.224, 0.225)$.
For SAM, we use the model-specific normalization with mean $(123.675, 116.28, 103.53)$ and standard deviation $(58.395, 57.12, 57.375)$.
All models are implemented in PyTorch and optimized using AdamW during fine-tuning.

